\documentclass[10pt,a4paper]{amsart}
\usepackage[margin=19mm]{geometry}
\usepackage{amsmath,amssymb,mathtools}
\usepackage[scr=rsfs,cal=euler]{mathalfa}
\usepackage{xfrac}
\usepackage[unicode,hidelinks,bookmarks=false]{hyperref}
\newcommand{\ie}{i.e.\ }
\newcommand{\cf}{cf.\ }
\newcommand{\resp}{respectively\ }
\newcommand{\Subsection}{\subsection}
\providecommand{\nobreakdash}{\nobreak\hbox{-}}
\setlength{\emergencystretch}{2em}
\multlinegap0pt
\usepackage{bm}
\usepackage{bbm}
\usepackage{dsfont}
\usepackage{xspace}
\usepackage{cancel}
\usepackage{mathtools}
\usepackage{amsthm}
\makeatletter
\@ifundefined{theorem}{\newtheorem{theorem}{Theorem}}{}
\@ifundefined{definition}{\newtheorem{definition}[theorem]{Definition}}{}
\makeatother
\newcommand{\Tau}{\mathrm{T}}
\title[Pseudo-differential operator associated with the linear cano]{Pseudo-differential operator associated with the linear canonical Dunkl transform}
\author{S. Umamaheswari, Sandeep Kumar Verma}
\date{Source: arXiv:2608.11535v1 (CC BY 4.0)}
\begin{document}
\maketitle

\noindent\emph{Source and licence.} Pseudo-differential operator associated with the linear canonical Dunkl transform. Version arXiv:2608.11535v1 (CC BY 4.0), distributed under the Creative Commons Attribution 4.0 International licence, as recorded in the arXiv licence metadata for this exact version and shown on the abstract page https://arxiv.org/abs/2608.11535v1. Full rights notice: licenses/arxiv-2608.11535-CC-BY-4.0.txt.\par\smallskip

\noindent\textit{Editorial scope.} Counted unit: the Theorem whose source label is \texttt{None} in the article (source0.tex lines 711--713, source label \texttt{None}), delivered together with the article's own complete original proof of it (source0.tex lines 714--794, one \texttt{proof} environment of 81 lines). No mathematics is written, repaired or re-derived: statement and proof are the article's own text line for line. Required context retained uncounted: C7-e:2.1 (lines 144--146); C7-D3.2 (lines 299--304). Every definition, display and label of those lines is retained unchanged. Omitted from this package and not redistributed: 1--143, 147--298, 305--710, 795--966. Nothing inside a retained excerpt is omitted or altered: every retained body is delivered exactly as the article's lines are written, so no omission receipt of any kind accompanies this package. Citations: the retained excerpts carry no citation command at all, so no bibliography entry of the article is transcribed and no reference is redistributed. Reference adaptation: none was needed and none was made; every reference inside the delivered unit resolves inside it, so no unresolved reference or citation is produced. Printed slips kept literal: no printed word, symbol, hypothesis or number of the counted statement or of its proof was altered. Layout and class: the article's own class is \texttt{amsproc} with packages including amsmath, amsthm, amscd, amsfonts, amssymb, graphicx, color, subcaption, hyperref, enumerate; the wrapper is the lane's pinned \texttt{amsart} editorial wrapper at 10pt on A4, which restates the theorem-like environments the retained text needs and adds the article's own macro definitions that the retained text needs (\texttt{\textbackslash Tau}), with extra wrapper packages (bm, bbm, dsfont, xspace, cancel, mathtools). The delivered numbering is the unit's own. Assets: no retained span contains a figure environment, a graphics-inclusion command, a PDF-page inclusion, a table or a TikZ picture; no image file is copied into this package, the package declares no asset dependency and no image file is redistributed with it. Licence: version arXiv:2608.11535v1 of this article is distributed under the Creative Commons Attribution 4.0 International licence, as recorded in the arXiv licence metadata for this exact version and shown on the abstract page https://arxiv.org/abs/2608.11535v1. A full rights notice is delivered at licenses/arxiv-2608.11535-CC-BY-4.0.txt. Characters: every retained excerpt is pure ASCII, so the delivered engine, XeLaTeX with its default Latin Modern text font, reports no missing character.
\begin{equation}\label{C7-e:2.1}
\mathcal{D}_k^M \left( \Lambda_{k, M^{-1}}^nf\right)(\lambda) = \left( \frac{i\lambda}{b}\right)^n \, \mathcal{D}_k^M(f)(\lambda),
\end{equation}

\begin{definition} \label{C7-D3.2}
Let  $a$ be a symbol in the class $S_0^m$. Then the pseudo differential operator $P_{a,M}$ on $ \mathcal{S}(\mathbb{R})$ is defined by
\begin{equation*}
P_{a,M}(f)(x) =  \int_{\mathbb{R}}  a(x,\lambda)\, \mathcal{D}_k^M(f)(\lambda)\, E_k^{M^{-1}}(x,\lambda)\,d\mu_k(\lambda), \quad \text{ for every} \quad f \in \mathcal{S}(\mathbb{R}).
\end{equation*}
\end{definition}

 \begin{theorem}  If $a\in S^m_0$, then the pseudo-differential operator 
 $P_{a,M}:L^2_k(\mathbb{R}) \to L^2_k(\mathbb{R})$ is a bounded operator.
\end{theorem}

\begin{proof}
We first prove this result for a symbol $a(x,\lambda)$, which is compactly supported with respect to the variable $x$. We can express $a(\cdot,\lambda)$ as
\begin{equation} \label{C7-e3.1}
 a(x,\lambda) = \int_{\mathbb{R}} \mathcal{D}_k^M(a(\cdot, \lambda))(\eta)\, E_k^{M^{-1}}(x,\eta)\, d\mu_{k,-b}(\eta).
\end{equation}
It is clear that  $a(\cdot, \lambda) \in \mathcal{S}(\mathbb{R})$. Therefore, using \eqref{C7-e:2.1} , we have 
\begin{equation*}
    \left(\frac{i\eta}{b}\right)^\alpha\, \mathcal{D}_k^M(a(\cdot, \lambda))(\eta) = \int_{\mathbb{R}} \Lambda^\alpha_{k,M^{-1}} a(x,\lambda)\,E_k^M(x,\eta)\,d\mu_{k,b}(x).
\end{equation*}
Also, we have 
\begin{eqnarray*}
 \mathop{\underset{\eta
 \in \mathbb{R}}{\text{sup}}}  \left|  \left(\frac{i\eta}{b}\right)^\alpha\, \mathcal{D}_k^M(a(\cdot, \lambda))(\eta) \right| \le C_\alpha\,\left(1+ \left|\frac{\eta}{b} \right|\right)^{-N}, \quad \forall \quad 
\alpha,  N \in \mathbb{N}.
\end{eqnarray*}
Using the Defintion \ref{C7-D3.2} and \eqref{C7-e3.1}, we have 
\begin{eqnarray*}
    P_{a,M}(f)(x) &=& \int_{\mathbb{R}} \int_{\mathbb{R}} E_k^{M^{-1}}(x,\lambda)\,E_k^{M^{-1}}(x,\eta)\,\mathcal{D}_k^M(a(\cdot,\lambda))(\eta)\,\mathcal{D}_k^M(f)(\lambda)\,\\ &&\times d\mu_{k,-b}(\eta)\, d\mu_k(\lambda)\\
    &=&\int_{\mathbb{R}}  (Sf)(\eta,x)\,d\mu_k(\eta),
\end{eqnarray*}
where
\begin{equation*}
(Sf)(\eta,x) = \int_{\mathbb{R}} E_k^{M^{-1}}(x,\lambda)\,E_k^{M^{-1}}(x,\eta)\,\mathcal{D}_k^M(a(\cdot,\lambda))(\eta)\,\mathcal{D}_k^M(f)(\lambda)\,d\mu_{k,-b}(\lambda).
\end{equation*}
Let us consider
\begin{eqnarray*}
    \int_{\mathbb{R}} |(Sf)(\eta,x)|^2\,d\mu_k(\eta) &\le&\int_{\mathbb{R}} |\mathcal{D}_k^M(a(\cdot,\lambda))(\eta)|^2\,|\mathcal{D}_k^M(f)(\lambda)|^2\,d\mu_{k,-b}(\lambda)\\
    &\le&  \mathop{\underset{\eta
 \in \mathbb{R}}{\text{sup}}}|\mathcal{D}_k^M(a(\cdot,\lambda))(\eta)|\,\|f\|_{L^2(\mathbb{R})}\\
 &\lesssim& \left(1+ \left|\frac{\eta}{b} \right|\right)^{-N}\,\|f\|_{L^2(\mathbb{R})}.
 \end{eqnarray*}
 Thus, we have
 \begin{eqnarray*}
 \|P_{a,M}f\|_{L_k^2(\mathbb{R})} &\le& \int_{\mathbb{R}}\|Sf(\eta,\cdot)\|_{L_k^2(\mathbb{R})}\,d\mu_k(\eta)\\
 &\lesssim& \int_{\mathbb{R}}\left(1+ \left|\frac{\eta}{b} \right|\right)^{-N}\,\|f\|_{L_k^2(\mathbb{R})}\,d\mu_k(\eta)\\
\|P_{a,M}f\|_{L_k^2(\mathbb{R})} &\lesssim&  \|f\|_{L_k^2(\mathbb{R})}.
 \end{eqnarray*}
 The result is proved for the compactly supported symbols. Next, we prove the result for the symbol $a$ which is not necessarily compactly supported.  For all $x_0 \in \mathbb{R}$ and $ N\ge0$,  there exist a constant $C_N(x_0
 )$ such that
 \begin{equation}\label{C7-e3.2}
\int_{|x-x_0| \le 1} |P_{a,M}(f)(x)|^2\,d\mu_k(x) \le  C_N(x_0
 )\, \int_{\mathbb{R}}\frac{|f(x)|^2}{(1+|x-x_0|^2)^N}\,d\mu_k(x).  
 \end{equation}
For the time being, we can assume that \eqref{C7-e3.2} is true and prove the boundedness of $P_{a,M}$.
\begin{eqnarray*}
    \int_{\mathbb{R}^2} \chi_{|x-x_0|\le 1}\, |P_{a,M}(f)(x)|^2\,d\mu_k(x)\,d\mu_k(x_0) &\lesssim& \int_{\mathbb{R}^2}  \frac{|f(x)|^2}{(1+|x-x_0|^2)^N}\,d\mu_k(x)\,d\mu_k(x_0), \end{eqnarray*}
    where $N >k+\frac{1}{2}$. Applying Fubini's theorem, we obtain that
\begin{eqnarray*}
l(\left[x-1,x+1\right])\,\int_{\mathbb{R}} |P_{a,M}(f)(x)|^2\,d\mu_k(x) &\lesssim& \int_{\mathbb{R}} |f(x)|^2\, d\mu_k(x)\\
&\lesssim& \int_{\mathbb{R}} |f(x)|^2\, d\mu_k(x).
\end{eqnarray*}
This completes the required proof. Now, let us 
  prove \eqref{C7-e3.2}. We start to prove this result for $x_0=0$ and then extend it to general $x_0\in \mathbb{R}$. The function $f=f_1+f_2$ can be written in terms of two functions $f_1$ and $f_2$ such that the support of $f_1$ contained in $\{x \in \mathbb{R}: |x|\le 3\}$ and support of $f_2$ contained in $\{x \in \mathbb{R}: |x|\ge 2\}$, respectively.  First, we prove the inequality  \eqref{C7-e3.2} for $f_1$.\\
Let $\gamma \in C_c^\infty(\mathbb{R})$ with $\gamma(x)=1$ on $|x|\le 1$. Then we define $\gamma(P_{a,M}) := (\gamma P)_{a,M}$ is a compactly supported pseudo-differential operator whose symbol is $\gamma(x)\,a(x,\lambda)$. We start with the estimation that
\begin{eqnarray}
\nonumber    \int_{-1}^{1} |P_{a,M}(f_1)(x)|^2\,d\mu_k(x) &=& \int_{\mathbb{R}} |
(\gamma P)_{a,M}(f_1)(x)|^2\, d\mu_k(x)\\
\nonumber  &\lesssim&    \int_{\mathbb{R}} |f_1(x)|^2\, d\mu_k(x) \\
  &\lesssim& \int_{-3}^3 |f(x)|^2\, d\mu_k(x).\label{C7-e:3.3}
\end{eqnarray}
Next, we shall prove the estimate for $f_2$. If $\{x\in \mathbb{R}:|x|\le 1\}$, then $x \notin \text{supp}\,f_2$. Therefore,  we have 
\begin{eqnarray*}
   | P_{a,M}(f_2)(x)| &\le& \int_{-2}^2 |f_2(y)|\, |\mathcal{K}^M(x,y)|\,d\mu_k(y)\\
    &\le&  \int_{\mathbb{R}}  |f(y)|\,|\Tau_x^M(K^M(x,\cdot))(y)|\,d\mu_k(y)
    \end{eqnarray*}
    Using Schwartz's inequality on the right-hand side of the integration, we get
    \begin{eqnarray*}
| P_{a,M}(f_2)(x)|    &\le& C\, \left(\int_{\mathbb{R}} \frac{|f(y)|^2}{(1+y^2)^N}\, d\mu_k(y)\right)^{\frac{1}{2}},
    \end{eqnarray*}
where 
\begin{equation*}
    C = \left(\int_{\mathbb{R}}(1+y^2)^N\,|\Tau_x^M(K^M(x,\cdot))(y)|^2\, d\mu_k(y)\right)^{\frac{1}{2}}.
\end{equation*}
Since $\Tau_x^M(K^M(x,\cdot))$ is a smooth function with suitable decay, therfore the above integral is converges.
Taking integration on both sides, we get
    \begin{eqnarray}\label{C7-e:3.4}
\int_{-1}^1  | P_{a,M}(f_2)(x)|^2\,d\mu_k(x)  &\lesssim&  \int_{\mathbb{R}} \frac{|f(y)|^2}{(1+|y|)^N}\, d\mu_k(y).
\end{eqnarray}
Combining \eqref{C7-e:3.3} and \eqref{C7-e:3.4} implies that \eqref{C7-e3.2}.
Now, we prove \eqref{C7-e3.2}   for an arbitrary $x_0\in \mathbb{R}$. We define $a_{x_0}(x,\lambda) = a(x-x_0, \lambda)$ and we see that the estimates of $P_{a,M}$ in $\{x\in \mathbb{R}:|x-x_0| \le 1\}$ is equivalent to  the same estimste of $P_{a_{x_0},M}$ in $\{x\in \mathbb{R}: |x|\le1\}$. This completes the proof of the theorem.
\end{proof} 

\end{document}
